\section{Introduction}
\begin{figure*}[h!]
    \centering
    \includegraphics[width=\textwidth]{pgs-flowchart.png}
    \caption{A simplified depiction of forensic lab use of probabilistic genotyping software to produce evidence in a criminal case. The goal of forensic DNA testing is to determine whether a person of interest (e.g., the defendant) could be the source of some or all of the DNA recovered from a crime scene. To this end, \textbf{(A)} a lab analyst first prepares two physical DNA samples: an evidence sample (e.g., a swabbing of a shoe collected as evidence) and a reference profile. \textbf{(B)} The analyst then runs the PGS with these DNA samples as inputs, along with other inputs chosen by the analyst. The software outputs a \textit{likelihood ratio} (LR)---a statistic that describes similarities observed between the evidence sample and reference profile. \textbf{(C)} Before the LR is presented as evidence at trial, a court may conduct an admissibility hearing to determine whether the evidence is sufficiently reliable to be presented in front of a jury. \textbf{(D)} If the judge deems the evidence admissible, the LR will be presented as evidence at trial.}.
    \label{fig:pgs-flowchart}
    \Description{Flow-chart depicting stages of the forensic DNA testing process and the introduction of the probabilistic genotyping software (PGS) likelihood ratio output as evidence in a criminal case. The flowchart has four stages: (A) DNA profile preparation, (B) Use of PGS to produce likelihood ratio, (C) A potential admissibility hearing, and (D) Presentation of LR as evidence at trial. The figure depicts four inputs to the probabilistic genotyping software: the crime scene DNA sample (i.e., evidence sample), the person of interest's DNA sample (i.e., reference profile), the forensic lab analyst's estimate of the number of contributors, and other analyst inputs. The software outputs a likelihood ratio (LR), which the figure then depicts as being presented in front of a judge. If the judge deems the LR admissible, the LR is shown as trial evidence in front of a jury. There is a box depicting the PGS. There is also a larger box encompassing both the PGS and lab analyst with a dotted line depicting the broader PGS system that takes physical DNA samples as input and outputs a LR.}
\end{figure*}

Over the past decade, researchers and practitioners have developed an ecosystem of tools---including guidelines, frameworks, and software toolkits---to promote fairness, transparency, and accountability throughout the AI lifecycle~\cite{madaio2020co, holstein2019improving, lee2021landscape, raji2020closing, mitchell2019model, abebe2022adversarial}. Amidst growing calls for algorithm audits, recent work highlights that such tools can be crucial for ensuring these evaluations result in meaningful accountability~\cite{ojewale2025towards, abebe2022adversarial, lam2024framework, raji2022outsider, jin2024beyond}. However, tools for responsible AI may fail to achieve their intended outcomes when deployed in complex organizational settings and practices~\cite{wong2023seeing, berman2024scoping, deng2022exploring, rakova2021responsible, madaio2020co, schor2024mind}. 


In response, a burgeoning body of work investigates how to design audit tooling that effectively helps audits create accountability~\cite{ojewale2025towards, madaio2020co, raji2020closing, holstein2019improving, lam2024framework, raji2022outsider, abebe2022adversarial}. As researchers and policymakers increasingly call for audit standards (i.e., agreed-upon best practices)~\cite{goodman2022algorithmic, costanza2022audits, lam2024framework, raji2022outsider, groves2024auditing, aistandardshub}, we ask: \textit{How can the design of an audit standard undermine its effectiveness?} We specifically examine how a standard’s formal requirements may enable audits that comply with the letter of the standard while failing to meet the goals envisioned by the standard's developers. We explore this question through the lens of \textbf{ASB 018\footnote{We study the first edition of ASB 018, which was developed and published in 2020 by the American Academy of Forensic Sciences Standards Board (ASB)~\cite{asb018}.}}, a standard for auditing \textbf{probabilistic genotyping software (PGS)} used by U.S. forensic laboratories to analyze DNA samples in criminal cases. ASB 018 provides a critical case study not only because it mirrors audit processes in other domains like hiring~\cite{raji2022outsider, lam2024framework, groves2024auditing, gerchick2025auditing}, healthcare~\cite{obermeyer2019dissecting, wu2021medical}, and social services~\cite{chouldechova2018case} but also because forensic labs already cite ASB 018 compliance to convince judges that PGS outputs are sufficiently reliable to be presented as evidence in criminal trials~\cite{2023anderson}. 


We investigate ASB 018's design and effectiveness through two research questions: \textit{(RQ1) What gaps exist between the practices ASB 018 envisions and the practices it enables?\footnote{We study forensic labs' audit practices by examining their representation in publicly available audit reports; see Section~\ref{sec:methods}.}}, and \textit{(RQ2) How does the design of ASB 018 facilitate these gaps?} In this work, we analyzed five publicly available PGS audit reports spanning various lab sizes, software versions, and years. We first assessed each audit report for compliance with ASB 018’s formal requirements. After determining that all five audits could be interpreted as compliant, we qualitatively coded ASB 018 and its accompanying factsheet to identify its desired outcomes across five audit stages~\cite{ojewale2025towards}: Audit Scope, Standards Identification, Performance Analysis, Post-Audit Judgment, and Audit Communication. We then qualitatively coded each audit report to identify practices that misalign with the standard’s goals (RQ1). Finally, we conducted a second round of coding on ASB 018, this time focusing on the language and organization of its requirements to identify aspects of its design that allow unsatisfactory audits to remain technically compliant (RQ2).

Our findings reveal that the analyzed audits are ASB 018 \textit{compliant, but unsatisfactory}: \textbf{audit practices consistently fall short of ASB 018’s envisioned goals across all five audit stages, yet still satisfy the standard's formal requirements}. For instance, while ASB 018 envisions testing the broader sociotechnical system, the audits primarily measure narrow software performance. Similarly, although the standard intends for audits to establish boundaries on software use based on failures observed during performance analysis, the audits avoid framing errors as failures and establish no such boundaries. Furthermore, despite the standard's emphasis on documentation that supports third-party evaluation, the reports omit crucial details needed to verify audit claims. \textbf{ASB 018 enables these gaps through design choices such as undefined concepts, requirements language that treats audit components as separate considerations to be addressed in isolation, and vague language calling for labs to \textit{``address''} and \textit{``consider''} that allow compliance with minimal to no testing.}


To conclude, we discuss how our study demonstrates that clear articulation of a tool’s goals---even from an external perspective---is essential for evaluating and improving audit infrastructure. Towards designing audit tools that more effectively deliver accountability, we caution that designing for compatibility with users' current perspectives and practices can undermine the tool's effectiveness. To help navigate this tension, we provide recommendations for those designing standards to (1) evaluate standards against fine-grained articulations of desired outcomes to identify where increased specificity is needed, using frameworks such as that by \citet{ojewale2025towards} to articulate intended outcomes for audit stages beyond evaluation; (2) design requirements that clearly specify actor responsibilities, minimum activities, and intermediary steps essential to achieving desired outcomes; and (3) co-design standards with stakeholders beyond intended users to encourage more robust definitions of desired outcomes that counterbalance users' desires for flexibility. 



